\documentclass[17pt]{extarticle}

\usepackage{extsizes}

\usepackage[paperheight=9.6cm,paperwidth=17cm,margin=0.7cm]{geometry}
\parindent=0mm \parskip=3.2mm

\usepackage{amsmath,amssymb,mathtools}
\usepackage{graphicx,xcolor,tikz}
\usepackage{enumerate,enumitem}

\setlist{leftmargin=*}
\usepackage[russian]{babel}

\definecolor{fon}{HTML}{0f1622}

\pagecolor{fon}
\color{white}

\usepackage{mathspec}

\setmainfont[
	Path = f/,
	BoldFont=rb.ttf,
	ItalicFont=ri.ttf,
	BoldItalicFont=rbi.ttf
		]{r.ttf}
\setsansfont[
	Path = f/,
	BoldFont=rbb.ttf,
	ItalicFont=rli.ttf,
	BoldItalicFont=rbbi.ttf
		]{rl.ttf}
		
\setmathfont(Digits)[Path = f/]{rm.ttf}
\setmathfont(Latin)[Path = f/]{rmi.ttf}
\setmathfont(Greek)[Path = f/, Uppercase]{rm.ttf}
\setmathfont(Greek)[Path = f/, Lowercase]{rmi.ttf}

\setmonofont[Path = f/]{rm.ttf}

\DeclarePairedDelimiter\lr{(}{)}
\newcommand{\divsby}{\mathrel{\rlap{.}\rlap{\raisebox{0.55ex}{.}}\raisebox{1.1ex}{.}}}
\renewcommand{\alpha}{α}
\newcommand{\euler}{φ}
\newcommand{\newslide}[1]{\newpage \begin{center} \large #1 \end{center}}
\newcommand{\tr}[1]{\textcolor{red}{#1}}
\newcommand{\vnode}{node[circle,fill=white,fill opacity=1,inner sep=0pt,minimum size=1.7mm]{}}

\begin{document}

\ \\ [1cm]

\begin{center} \Large «Математика НОН-СТОП»:\medskip \\ 2021 И 2022 \end{center}

\newslide{Олимпиаду проводит Фонд «Время Науки»}

\begin{center}
%	\includegraphics[height=2.6cm]{fig/fpg}\quad
	\includegraphics[height=3.2cm]{fig/vn}
\end{center}

\newslide{«Математика НОН-СТОП» — 2021}

\begin{itemize}
	\item Более 2500 проверенных работ,
	\item Половина участников получила\\ дипломы и п/о,
	\item 24 площадки, Москва и Калининград, {\it Элиста,}
	\item \( 28 \cdot 3\) базовых задач,\\ распределённых по вариантам
\end{itemize}

\newslide{Площадки олимпиады}

\begin{itemize}
	\item В 2022, если площадка «только для своих»,\\
	регистрация всё равно будет\\
	через {\it rs.mathnonstop.ru}, формат с таблицами\\
	замедлил формирование списков.
	\item Сохранятся ли ограничения на рассадку?\\ Надо ли формировать её заранее?
\end{itemize}

\newslide{Организация олимпиады}

\begin{itemize}
	\item {\it Концепция взаимодействия площадок:}\\
	фиксирует инструкции для организаторов,\\
	слова волонтёров перед олимпиадой,\\
	порядок регистрации и проверки,\\
	правила оценивания работ и контроля.
	\item В 2022~— наблюдатель от фонда «ВН»\\
	на каждой площадке.
\end{itemize}

\newslide{Организация проверки}

\begin{itemize}
	\item Коллектив одной школы отвечает за\\
	один вариант олимпиады целиком.
	\item Регионы проверяют у себя и могут\\
	сами вручать собственные призы.
	\item Все профильные варианты\\
	проверяют составители задач.
\end{itemize}

\newslide{Новые приёмы в задачах}

Каждый год мы пишем задачи\\
для олимпиады с нуля, и каждый год находим\\
для себя новые приёмы в задачах,\\
обращаем внимание на что-то новое.

\newslide{4 класс, 4B}

Вдохновлена олимпиадой «только на ответ» ЯКласс.\\ Задача про «пропорцию с хвостом».

Если корабль длиной 250 м идёт на полной скорости,\\ то нужно целиться на 250 м перед носом корабля.\\ Корабль идёт с $\frac{1}{3}$ максимальной скорости.\\ В какую точку нужно целиться?
	\[ \lr*{\lr*{250 + 125} / 3} - 125 = 0. \]

\newslide{5 класс, 6}

\begin{center}
\tikz[scale=0.6]{
	\draw[very thick] (0,0) -- (0,3) -- (1,3); \draw (0.5,-0.7) node {\large $1$};
	\draw[very thick] (2,0) -- (2,3); \draw[very thick] (2,2) -- (3,2); \draw (2.5,-0.7) node {\large $2$};
	\draw[very thick] (4,0) -- (4,3) -- (5,2); \draw (4.5,-0.7) node {\large $3$};
	\draw[very thick] (6,0) -- (6,3); \draw[very thick] (6,2) -- (7,3); \draw (6.5,-0.7) node {\large $4$};
	\draw[very thick] (8,0) -- (8,3) -- (9,3) -- (8,2); \draw (8.5,-0.7) node {\large $5$};
	\draw[very thick] (10,0) -- (10,3);\draw[very thick] (11,2) -- (11,3); \draw (10.5,-0.7) node {\large $6$};
	\draw[very thick] (12,0) -- (12,3) -- (13,3) -- (13,2); \draw (12.5,-0.7) node {\large $7$};
	\draw[very thick] (14,0) -- (14,3); \draw[very thick] (14,2) -- (15,2) -- (15,3); \draw (14.5,-0.7) node {\large $8$};
	\draw[very thick] (16,0) -- (16,3) -- (17,3) -- (17,2) -- (16,2); \draw (16.5,-0.7) node {\large $9$};
} 
\end{center}\vspace{-4mm}

Любое число от 1 до 9999 можно однозначно\\
восстановить по его цистерцианской записи.

\newslide{6 класс, 1A} \vspace{-0.4cm}

Давно интересовавший вопрос:

\begin{center} \includegraphics[width=0.55\paperwidth]{fig/cards-dark} \end{center} \vspace{-0.4cm}

Окажутся ли две карты, изначально бывшие рядом,\\
снова рядом после двух таких перемешиваний?

\newslide{7 класс, 3C}

Задача, которую никто не решит:\\
Придумать число, цифры которого при умножении\\
на \( 1\ldots 6 \) остаются теми же,\\
меняется лишь порядок.

{\small
\begin{center}\begin{tabular}{llll}
   $142857 \cdot 1 = 142857$ & $10^0 \equiv 1$ & $142857 \cdot 4 = 571428$ & $10^4 \equiv 4$ \\
   $142857 \cdot 2 = 285714$ & $10^2 \equiv 2$ & $142857 \cdot 4 = 714285$ & $10^5 \equiv 5$ \\
   $142857 \cdot 3 = 428571$ & $10^1 \equiv 3$ & $142857 \cdot 4 = 857142$ & $10^3 \equiv 6$ \\
\end{tabular}\end{center}}

\newslide{8 класс, 2A}

Аккуратный подсчёт: \begin{center} \tikz[scale=0.9]{
	     \draw[thick] (-2,0) \vnode -- (-1,0) \vnode -- (0,0) \vnode --
	          (0,1) \vnode -- (-1,1) \vnode -- (-2,1) \vnode;
	     \draw[thick] (0,1) -- (1,1) \vnode -- (2,1) \vnode --
	          (3,1) \vnode -- (4,1) \vnode;
	     \draw[thick] (0,0) -- (1,1);
	} \end{center} \vspace{-8mm}

\[C_3^2 + C_3^2 + C_4^2 +
	2 \cdot 3 \cdot 3 + 2 \cdot 3 \cdot 4 + 2 \cdot 3 \cdot 4.\] \vspace{-8mm}

Формулой, как обычно, посчитать быстрее\\ и надёжнее, чем перечисляя все пути.


\newslide{Методические пособия}

\begin{itemize}
	\item Решения «Математика НОН-СТОП — 2020»:\\ выдаётся призёрам и победителям 2021 года.
	\item Решения «Математика НОН-СТОП — 2021»:\\ публикация зимой 2021–2022.
	\item Задачи Санкт-Петербургских турниров юных\\
	математиков с комментариями авторов:\\
	опубликована, распространяется по школам.
\end{itemize}

\newslide{Другие мероприятия Фонда «Время Науки»}

\begin{itemize}
	\item Санкт-Петербургский Турнир \\
		юных математиков: \\
		старшая и младшая лиги
	\item Балтийский научно- \\
		инженерный конкурс
\end{itemize}

\newpage \begin{center} \Large Спасибо за внимание! \end{center}

\begin{itemize}
	\item Статистика и достижения 2021
	\item Организация 2022 $^*$
	\item Особенности задач 2021
\end{itemize}

\end{document}
